% FILE: 06_contribution_to_thesis.tex
% Project: adversarial
% Thesis Role: adversarial-testing-and-hostile-assumption-verification

\section{Contribution to the Dissertation}
\label{sec:adversarial-contribution}

\subsection{Contribution to Prediction vs.~Coordination}
\label{sec:adversarial-contribution-prediction}

The central thesis argument distinguishes prediction from coordination. Prediction produces a
probability estimate about a future state; coordination produces a cryptographically certified
commitment that a specified state transition occurred. The adversarial corpus demonstrates
that prediction without coordination is adversarially exploitable: a prediction model trained
on unverified event logs can be manipulated by an adversary who controls the log, because the
model has no mechanism to distinguish genuine from forged process evidence.

The simulation engine makes this concrete. Scenario 3 demonstrates that an adversary who
manufactures traces under a forged key produces event sequences that are structurally
indistinguishable from genuine traces when viewed as raw timestamps and activity names. Only
the HMAC-SHA256 signature check against the genuine system key detects the forgery. A
prediction model that operates on raw event sequences---without a prior coordination layer
enforcing signature authenticity---would train on the forged trace as if it were genuine
process behavior, corrupting all downstream predictions.

The adversarial corpus therefore provides the negative proof for the prediction-versus-coordination
distinction: coordination is not an enhancement to process intelligence; it is the prerequisite
that makes process intelligence non-exploitable.

\subsection{Contribution to Process Evidence}
\label{sec:adversarial-contribution-evidence}

The adversarial corpus contributes three specific findings to the process evidence
literature. First, the PROV-O gap analysis (\texttt{standards/ocel-adversarial-gravity.md})
establishes that XES and OCEL 2.0 are descriptive serializations, not cryptographically sound
operational ledgers. Neither standard enforces \texttt{wasGeneratedBy}, \texttt{used}, or
\texttt{wasAttributedTo} as cryptographic primitives, making them structurally vulnerable to
the five attack vectors cataloged in the document: XES flat-earth lossy ETL, OCEL 2.0
phantom record injection, object-to-event desynchronization, time-of-observation versus
time-of-occurrence manipulation, and lack of actor identity binding.

Second, the adversarial canon (\texttt{sources/papers/adversarial-canon.md}) establishes
that classical discovery algorithms (Alpha, Heuristics Miner, Inductive Miner) structurally
fail under Sybil attack. The canonical algorithms rely on temporal proximity and case
identifiers for correlation; a Sybil attacker can forge thousands of interleaved case
identifiers and events, causing state-space explosion and invalid spaghetti model discovery.
The absence of BLAKE3 cryptographic receipts means there is no provenance filter for
adversarial events.

Third, the Chatman Constraint and the three structural laws
(\texttt{sources/wasm4pm-compat/adversarial-type-law.md}) contribute a formal type-law
framework for witness lattice integrity: strict injectivity prevents state-space aliasing,
exhaustive typestate closure eliminates reachable undefined states, and cryptographic isolation
via domain separators prevents cross-epoch witness contamination.

\subsection{Contribution to Manufacturing Architecture}
\label{sec:adversarial-contribution-manufacturing}

The adversarial corpus defines the manufacturing pipeline's rejection boundary. The
CodeManufactory manufacturing pipeline cannot claim that a manufacturing receipt is valid
unless the ingestion boundary has verified the input against all five detection rules. The
adversarial corpus is the specification of those rules and the executable test suite that
verifies them.

The seven challenge categories in \texttt{RED\_TEAM\_SCOPE.md} also contribute to
manufacturing architecture by defining what constitutes an admissible board artifact. A claim
that survives all seven challenge categories---unsupported doctrine, inflated type coverage,
evidence-less board claims, observation-only lifecycle phases, evidence-less comparison
claims, loss-policy-free crosswalks, and diligence-path-free M\&A claims---is a manufacturing
receipt. A claim that does not survive is a research claim, appropriately marked and
sequestered from downstream manufacturing workflows.

The gap remediation sequencing contributed by Finding 010
(\texttt{adversarial/RED\_TEAM\_FINDINGS\_002.md}) defines the ordered authorization sequence
for downstream workflow agents: GAP\_007 first (independent, low effort), then GAP\_001
(critical path unlock), then GAP\_002, then parallel execution of GAP\_003, GAP\_004, and
GAP\_008, then parallel execution of GAP\_005 and GAP\_006. This sequencing is a direct
output of the adversarial challenge to board claims without evidence paths.

\subsection{Contribution to Capital-Flow Infrastructure}
\label{sec:adversarial-contribution-capital}

The adversarial corpus establishes the evidentiary requirements for capital-flow
infrastructure. The OCEL adversarial gravity analysis concludes that deploying XES or OCEL
without a cryptographic provenance wrapper in high-stakes environments is architectural
negligence and that M\&A boards cannot rely on process models derived from highly mutable
ledgers. The Reality Receipt mechanism---wasm4pm acting as an interception proxy that
encapsulates raw event data in a signed cryptographic envelope---is the minimum viable
infrastructure for board-admissible process evidence.

Finding 002 (\texttt{adversarial/RED\_TEAM\_FINDINGS\_001.md}) establishes the important
and self-correcting result that M\&A documents in \texttt{ma/} correctly identify themselves
as research claims, not manufacturing receipts. The distinction matters for capital-flow
infrastructure because a capital decision made on the basis of a research claim rather than a
manufacturing receipt is a governance failure. The adversarial corpus enforces this
distinction by requiring that every board projection trace to a manufacturing receipt through
the path: \emph{target OCEL log} $\rightarrow$ \emph{conformance report} $\rightarrow$
\emph{board projection}.
